\documentclass[conference]{IEEEtran}
\usepackage{iftex}
\ifPDFTeX
  \usepackage[utf8]{inputenc}
  \usepackage[T1]{fontenc}
\fi
\usepackage{amsmath,amssymb,bm,mathtools}
\usepackage{newtxtext,newtxmath}
\usepackage{textcomp}
\usepackage{enumitem}
\usepackage{booktabs}
\usepackage{siunitx}
\usepackage[breaklinks=true,hidelinks]{hyperref}
\usepackage[nameinlink,capitalise,noabbrev]{cleveref}
\usepackage[final]{microtype}
\interdisplaylinepenalty=2500
\allowdisplaybreaks
\DeclareMathOperator{\clamp}{clamp}
\DeclareMathOperator{\atanTwo}{atan2}
\newcommand{\vct}[1]{\bm{#1}}
\newcommand{\R}{\mathbb{R}}
\newcommand{\dd}{\,\mathrm{d}}
\title{Adaptive Vertical PII with Mahony Attitude and Wave-Frequency Scheduling\\(implementation note)}
\author{\IEEEauthorblockN{OpenAI Codex} \IEEEauthorblockA{Repository documentation draft}}
\begin{document}
\maketitle
\begin{abstract}
This note documents the method implemented by \texttt{w3d\_sim\_adaptive\_pii\_mahony.cpp} and the supporting headers \texttt{AdaptiveVerticalPIIMahony.h}, \texttt{AdaptiveVerticalPII.h}, \texttt{WaveFrequencyTracker.h}, \texttt{VerticalPIIObserver.h}, and \texttt{Mahony\_AHRS.h}. The pipeline estimates attitude with a Mahony complementary filter, rotates body-frame specific force into a Z-up world frame, removes gravity to obtain vertical inertial acceleration, and drives an adaptive vertical displacement observer whose gains are scheduled from an internal dominant-wave-frequency tracker and an acceleration-variance estimate.
\end{abstract}
\section{Signal flow from the simulation driver}
The simulation entry point instantiates \texttt{AdaptiveVerticalPIIMahony<float,true>} as the heave filter. Gyroscope and accelerometer samples arrive in the runner's NED convention and are remapped back to the original simulator body Z-up convention before the observer is updated. Magnetometer samples are retained in the runner convention and reused with zero-order hold on every IMU tick whenever heading aiding is enabled. The reported displacement, velocity, and filtered acceleration are the scalar vertical states produced by the wrapped adaptive observer, while roll, pitch, and yaw come from the Mahony attitude state.

At sample time $k$ with period $\Delta t_k$, the implemented cascade is
\begin{align}
\vct{\omega}_{b,k} &\xrightarrow{\text{Mahony}} \hat{q}_{wb,k}, \\
\vct{f}_{b,k} &\xrightarrow{\hat{q}_{bw,k}} \hat{\vct{f}}_{w,k}, \\
\hat{a}_{z,k} &= \hat{f}_{w,z,k} - g, \\
\hat{a}_{z,k} &\xrightarrow{\text{Adaptive Vertical PII}} (\hat{p}_k,\hat{v}_k,\hat{S}_k,\hat{b}_k),
\end{align}
where $\hat{q}_{bw,k}$ is the body-to-world quaternion, $\hat{p}$ is vertical displacement, $\hat{v}$ is vertical velocity, and $\hat{S}$ is the integral of displacement used to suppress long-term drift.

\section{Mahony attitude subsystem}
\subsection{Continuous-time idea}
The Mahony AHRS estimates the quaternion of sensor/body frame relative to the auxiliary world frame. The implementation stores the world-to-body quaternion
\begin{equation}
\hat{q}_{wb} = [q_0\; q_1\; q_2\; q_3]^\top,
\end{equation}
and uses normalized accelerometer and optionally normalized magnetometer measurements as direction observations. With accelerometer only, the gravity-direction innovation is the cross product between measured and predicted unit gravity directions. Denoting the innovation by $\vct{e}=[e_x\;e_y\;e_z]^\top$, the corrected gyro is
\begin{equation}
\vct{\omega}_c = \vct{\omega}_m + \vct{i}_{\omega} + 2K_p\,\vct{e},
\end{equation}
with integral feedback
\begin{equation}
\dot{\vct{i}}_{\omega} = 2K_i\,\vct{e}.
\end{equation}
The quaternion kinematics are then
\begin{equation}
\dot{\hat{q}}_{wb} = \tfrac12\,\hat{q}_{wb}\otimes [0\;\vct{\omega}_c],
\end{equation}
followed by renormalization. In code this is implemented as a first-order Euler step in quaternion coordinates and explicit normalization each sample.

\subsection{Innovation used by the code}
For the IMU-only branch, the accelerometer is normalized and the predicted body-frame direction of world $+Z$ is extracted from the current quaternion as
\begin{equation}
\hat{\vct{g}}_b =
\begin{bmatrix}
2(q_1 q_3 - q_0 q_2)\\
2(q_0 q_1 + q_2 q_3)\\
q_0^2 - q_1^2 - q_2^2 + q_3^2
\end{bmatrix}.
\end{equation}
If $\hat{\vct{a}}_b$ is the normalized accelerometer direction, the innovation is the cross product
\begin{equation}
\vct{e} = \hat{\vct{a}}_b \times \hat{\vct{g}}_b.
\end{equation}
With magnetometer aiding, the same gravity term is augmented by a magnetic-field innovation constructed from the estimated Earth field in the body frame.

\subsection{Sea-state scheduling of Mahony gains}
The wrapper also adapts the Mahony gains according to how trustworthy the accelerometer looks as a gravity cue. It computes three scalars:
\begin{align}
\varepsilon_{\mathrm{norm}} &= \frac{\left|\lVert \vct{a}_b\rVert - g\right|}{g},\\
\varepsilon_{\mathrm{innov}} &= \lVert \hat{\vct{a}}_b \times \hat{\vct{g}}_b \rVert,\\
\sigma_a &= \text{vertical acceleration RMS from the core observer wrapper}.
\end{align}
Each term is turned into a trust score with one-pole smoothing and clamping; qualitatively, near-static norm agreement, small innovation, and low vertical dynamics imply higher trust. The final trust $\gamma\in[\gamma_{\min},1]$ interpolates between calm-water and rough-water gains,
\begin{align}
2K_p(\gamma) &= 2K_{p,\mathrm{rough}} + \gamma\bigl(2K_{p,\mathrm{calm}}-2K_{p,\mathrm{rough}}\bigr),\\
2K_i(\gamma) &= 2K_{i,\mathrm{rough}} + \gamma\bigl(2K_{i,\mathrm{calm}}-2K_{i,\mathrm{rough}}\bigr).
\end{align}
Hence strong wave-induced specific force weakens attitude correction so the Mahony filter does not chase non-gravitational acceleration.

\section{Body-to-world acceleration conversion}
The Mahony wrapper stores $\hat{q}_{wb}$, but the vertical observer needs world-frame inertial acceleration. Therefore the implementation conjugates the quaternion to obtain body-to-world rotation $\hat{q}_{bw}=\hat{q}_{wb}^{\ast}$ and applies the standard quaternion-vector formula
\begin{equation}
\hat{\vct{f}}_w = \mathcal{R}(\hat{q}_{bw})\,\vct{f}_b.
\end{equation}
The code uses the equivalent cross-product form
\begin{align}
\vct{t} &= 2\,\vct{q}_v \times \vct{f}_b,\\
\hat{\vct{f}}_w &= \vct{f}_b + q_w\vct{t} + \vct{q}_v \times \vct{t},
\end{align}
where $\hat{q}_{bw}=[q_w\;\vct{q}_v]$. Since the world frame is explicitly treated as Z-up, the vertical inertial acceleration passed downstream is
\begin{equation}
\hat{a}_z = \hat{f}_{w,z} - g.
\end{equation}
At rest, $\hat{f}_{w,z}\approx g$, so $\hat{a}_z\approx 0$.

\section{Vertical PII observer model}
\subsection{State equations}
The heave estimator itself is implemented in \texttt{VerticalPIIObserver}. Its states are filtered acceleration $a_f$, velocity $v$, displacement $p$, and the integral of displacement $S$. When the bias channel is enabled, it also keeps a slow trend state $d$ and an estimated acceleration bias $b$. The continuous-time model documented in the header is
\begin{align}
\dot{a}_f &= \frac{a_z-a_f}{\tau_a}, \\
\dot{S} &= p, \\
\dot{p} &= v, \\
\dot{v} &= a_f - b - k_v v - k_p p - k_s S, \\
\dot{d} &= \frac{p-d}{\tau_d}, \\
\dot{b} &= k_b d - \lambda_b b.
\end{align}
The first four equations define a third-order PII stabilization loop around the double-integrator $\ddot{p}=a_f-b$. The last two equations are a very slow bias estimator that interprets sustained low-frequency displacement trend as evidence of accelerometer bias.

\subsection{Repeated-pole design}
Instead of tuning $k_v$, $k_p$, and $k_s$ independently, the implementation parameterizes the PII core with a repeated real pole at $-r$:
\begin{equation}
(s+r)^3 = s^3 + 3rs^2 + 3r^2s + r^3.
\end{equation}
Matching coefficients yields
\begin{equation}
 k_v = 3r, \qquad k_p = 3r^2, \qquad k_s = r^3.
\end{equation}
This matters because the closed-loop shape stays overdamped and monotone as $r$ changes. The parameter $r$ simply scales the observer bandwidth without changing the all-real-pole structure.

\subsection{Discrete implementation}
For sample interval $\Delta t$, the code uses exact one-pole smoothing for the low-pass blocks through
\begin{equation}
\alpha(\Delta t,\tau)=\frac{\Delta t}{\tau+\Delta t},
\end{equation}
so the filtered acceleration update is
\begin{equation}
 a_{f,k+1} = a_{f,k} + \alpha(\Delta t,\tau_a)\,(a_{z,k}-a_{f,k}).
\end{equation}
The bias trend channel uses the same smoothing for $d$, then Euler integration for $b$:
\begin{align}
 d_{k+1} &= d_k + \alpha(\Delta t,\tau_d)(p_k-d_k),\\
 b_{k+1} &= b_k + \Delta t\,(k_b d_{k+1}-\lambda_b b_k).
\end{align}
The PII core is updated semi-implicitly,
\begin{align}
\dot{v}_k &= a_{f,k+1} - b_{k+1} - k_v v_k - k_p p_k - k_s S_k,\\
 v_{k+1} &= v_k + \Delta t\,\dot{v}_k,\\
 p_{k+1} &= p_k + \Delta t\,v_{k+1},\\
 S_{k+1} &= S_k + \Delta t\,p_{k+1}.
\end{align}
The state clamps in the implementation act as safety saturations rather than part of the nominal model.

\section{Adaptive scheduling in \texttt{AdaptiveVerticalPII}}
\subsection{Acceleration variance estimate}
The wrapper around \texttt{VerticalPIIObserver} maintains running mean and variance estimates of the vertical inertial acceleration:
\begin{align}
\mu_{k+1} &= \mu_k + \alpha_\mu(a_{z,k}-\mu_k),\\
\nu_{k+1} &= \nu_k + \alpha_\nu\bigl((a_{z,k}-\mu_{k+1})^2-\nu_k\bigr),\\
\sigma_{a,k+1} &= \sqrt{\max(\nu_{k+1},\sigma_{\min}^2)}.
\end{align}
This $\sigma_a$ is the sea-state intensity indicator used in both the vertical observer scheduler and the Mahony trust logic.

\subsection{Scheduling law}
Whenever auto scheduling is enabled and the scheduler timer expires, \texttt{AdaptiveVerticalPII} obtains a dominant acceleration frequency estimate $f_{\mathrm{sched}}$ from \texttt{WaveFrequencyTracker} and calls
\begin{equation}
\texttt{observer.update\_adaptation}(f_{\mathrm{sched}},\sigma_a,c,\Delta t_{\mathrm{sched}}).
\end{equation}
Inside \texttt{VerticalPIIObserver}, the tracker inputs are first smoothed:
\begin{align}
 f_f^+ &= f_f + \alpha_{\mathrm{in}}(f_{\mathrm{disp}}-f_f),\\
 \sigma_f^+ &= \sigma_f + \alpha_{\mathrm{in}}(\sigma_a-\sigma_f).
\end{align}
Then normalized ratios relative to reference sea-state values are formed,
\begin{equation}
\rho_f = \clamp\!\left(\frac{f_f}{f_{\mathrm{ref}}},0.25,4\right),\qquad
\rho_\sigma = \clamp\!\left(\frac{\sigma_{\mathrm{ref}}}{\sigma_f},0.25,4\right).
\end{equation}
The commanded parameters are power-law schedules,
\begin{align}
 r_{\mathrm{cmd}} &= r_0\,\rho_f^{\eta_r}\,\rho_\sigma^{\zeta_r},\\
 \tau_{a,\mathrm{cmd}} &= \tau_{a,0}\,\rho_f^{\eta_a}\,\rho_\sigma^{\zeta_a},\\
 \tau_{d,\mathrm{cmd}} &= \tau_{d,0}\,\rho_f^{\eta_d}\,\rho_\sigma^{\zeta_d},\\
 k_{b,\mathrm{cmd}} &= k_{b,0}\,\rho_f^{\eta_b}\,\rho_\sigma^{\zeta_b},
\end{align}
with hard parameter bounds after each equation. Finally the active parameters move toward the commands via one-pole smoothing,
\begin{equation}
\theta^+ = \theta + \alpha_p(\theta_{\mathrm{cmd}}-\theta),
\end{equation}
for each scheduled parameter $\theta\in\{r,\tau_a,\tau_d,k_b\}$.

The signs of the exponents encode the intended control effect. Higher dominant wave frequency generally permits a faster observer, so the implementation uses positive frequency exponent for $r$ and negative one for $\tau_a$. Higher acceleration variance indicates rougher seas, so it tends to reduce aggressiveness by shrinking $r$ and $k_b$ while increasing effective filtering through $\tau_a$ and $\tau_d$ according to the configured exponents.

\section{Wave frequency tracker model}
\subsection{Front-end preprocessing}
\texttt{WaveFrequencyTracker} is a single-tone dominant-frequency tracker operating on vertical inertial acceleration. Let the input be $x(t)$. The code first removes quasi-DC drift using a one-pole high-pass formed as a low-pass state subtraction,
\begin{align}
 x_{\mathrm{dc}}^+ &= \operatorname{LP}_{f_{\mathrm{hp}}}(x),\\
 x_{\mathrm{hp}} &= x - x_{\mathrm{dc}},
\end{align}
and then applies another one-pole low-pass to produce a broad band-passed signal $y$. Its RMS energy is estimated with another smoothed square law.

\subsection{Synchronous demodulation and PLL/FLL loop}
An internal oscillator with phase $\phi$ and angular frequency $\omega$ is maintained. The band-passed signal is synchronously demodulated,
\begin{equation}
I_m = y\cos\phi, \qquad Q_m = -y\sin\phi,
\end{equation}
then low-pass filtered to obtain $I$ and $Q$. For a narrowband input $y\approx A\cos(\phi+\delta)$, the low-pass outputs satisfy
\begin{equation}
 I\approx \frac{A}{2}\cos\delta, \qquad Q\approx \frac{A}{2}\sin\delta,
\end{equation}
so the phase error estimate is
\begin{equation}
 e_\phi = \atanTwo(Q,I).
\end{equation}
The loop then behaves like a type-II PLL/FLL,
\begin{align}
 \dot{\omega} &= K_i\,g_c\,e_\phi,\\
 \dot{\phi} &= \omega + K_p\,g_c\,e_\phi,
\end{align}
where
\begin{equation}
K_p = 2\zeta\omega_n,\qquad K_i = \omega_n^2,\qquad \omega_n = 2\pi f_{\mathrm{bw}}.
\end{equation}
The gain factor $g_c$ is derived from confidence and reduced further when the tracker is not locked. Frequency slew-rate limiting, band clipping, and slow recentering toward the nominal frequency protect the loop from divergence.

\subsection{Confidence and coarse assist}
The tracker computes amplitude estimate $\hat{A}=2\sqrt{I^2+Q^2}$, coherence
\begin{equation}
\chi = \clamp\!\left(\frac{\hat{A}}{\sqrt{2}\,\mathrm{RMS}(y)},0,1\right),
\end{equation}
and confidence as a smoothed product of coherence and RMS score. A hysteretic lock detector uses RMS and coherence thresholds.

A second, slower acquisition aid estimates a coarse period from alternating threshold crossings of the band-passed signal. If thresholded negative-to-positive events occur one period apart, then
\begin{equation}
\hat{f}_{\mathrm{coarse}} = \frac{1}{T_{\mathrm{event}}},
\end{equation}
smoothed over time. When lock is weak, the PLL frequency is gently pulled toward this coarse estimate. \texttt{AdaptiveVerticalPII} may also blend the coarse and PLL estimates before scheduling the observer.

\section{Overall interpretation of the coupled method}
The complete algorithm is a cascaded observer for heave in rough seas:
\begin{enumerate}[leftmargin=*]
\item Mahony AHRS estimates orientation from gyro integration corrected by gravity and optional magnetometer references.
\item The body specific force is rotated into a Z-up world frame and converted into vertical inertial acceleration by subtracting gravity.
\item A narrowband wave-frequency tracker estimates the dominant sea-state frequency and a confidence score from the same vertical acceleration signal.
\item A vertical PII observer integrates the vertical acceleration to displacement while using repeated-pole damping, displacement integral feedback, and a very slow bias channel to prevent unbounded drift.
\item The observer bandwidth and bias-learning rates are continuously scheduled from estimated wave frequency and signal variance, so calm seas allow tighter tracking while rough, high-variance conditions slow the loop and reduce sensitivity to spurious bias inference.
\end{enumerate}

This design avoids a full stochastic state-space model yet still introduces physically meaningful structure: attitude resolves gravity, the PLL extracts dominant sea period, and the adaptive PII loop shapes the heave dynamics to that estimated sea state.

\section{Notes specific to the simulation wrapper}
The simulation adapter in \texttt{w3d\_sim\_adaptive\_pii\_mahony.cpp} selects concrete gains and adaptation exponents. In particular, it enables automatic scheduling from the internal acceleration-frequency tracker, configures confidence floors so adaptation does not stall before lock, and applies separate calm/rough gain limits to Mahony. Yaw is reported relative to the first valid heading so each run starts from zero heading offset. The filter snapshot exported by the adapter exposes both raw and scheduled sea-state variables, including active $r$, $\tau_a$, $\tau_d$, $k_b$, estimated acceleration variance, estimated frequency, and implied displacement/velocity scales $\sigma_a/\omega^2$ and $\sigma_a/\omega$.

\end{document}
